% !TeX program = pdflatex
% =============================================================================
% Verlaufsplanung: Kondensatoren (Doppellektion) -- Lektionsablauf-Tabelle
% UZH Praktikum I -- Sara Romer Urban, Kantonsschule im Lee, HS2026
% Klassen: 4fh (regulaer, 25.08.2026) + 4cdeg (SP4, 28.08.2026)
%
% Vollstaendig ueberarbeitet am 2026-08-27 anhand von
% kondensatoren-lektionsplan.md (selber Ordner) -- ersetzt den vorherigen
% Ablauf (Coulombkraft-Herleitung, Sortieraufgabe vor Theorie 1), der im
% Arbeitsblatt/Aufgabenblatt nicht mehr existiert. Aktueller Aufbau: staendiger
% Wechsel Aufgabenblatt (PhET-/Excel-Lernaufgabe, Partnerarbeit) <-> Arbeitsblatt
% (Theorie, Plenumsboxen, Demoversuch, Dielektrikum-Herleitung).
%
% Layout an Zeitplan_Vorlage.tex angelehnt (longtable, alternierende
% Zeilenfarben). Spalten: Inhalt, Zeit, Lehrperson (Taetigkeit), SuS
% (Taetigkeit), Material, Sozialform. Die Impulsfragen der Lehrperson und
% die erwartbaren Schuelerfragen stehen (mit Antwort + Erklaerung) separat
% in kondensatoren-fragenkatalog.tex/.pdf.
% =============================================================================
\documentclass[11pt,a4paper]{article}

\usepackage[landscape,margin=1.6cm]{geometry}
\usepackage[ngerman]{babel}
\usepackage[T1]{fontenc}
\usepackage[utf8]{inputenc}
\usepackage{lmodern}
\usepackage{array}
\usepackage{longtable}
\usepackage{booktabs}
\usepackage[table]{xcolor}
\usepackage{enumitem}
\usepackage{amsmath}
\usepackage{fancyhdr}
\usepackage{hyperref}

\pagestyle{fancy}
\fancyhf{}
\lhead{Physik -- Kondensatoren}
\rhead{Lektionsablauf}
\cfoot{\thepage}

\setlength{\parindent}{0pt}
\setlength{\parskip}{4pt}

\definecolor{headgray}{RGB}{230,230,230}
\definecolor{lightrow}{RGB}{245,245,245}

\hypersetup{
  colorlinks=true,
  urlcolor=blue,
  linkcolor=blue,
  pdftitle={Lektionsablauf: Kondensatoren},
  pdfauthor={Loris Keller}
}

\begin{document}

\begin{center}
    {\LARGE \textbf{Lektionsablauf: Kondensatoren}}\\[0.3cm]
    \large Elektrodynamik -- Doppellektion (2$\times$45 min) -- 4fh/4cdeg, HS2026
\end{center}

\vspace{0.2cm}
{\small Ständiger Wechsel zwischen \textbf{Aufgabenblatt 3} (PhET-/Excel-Lernaufgabe,
Partnerarbeit am Computer, Teil 1--4) und \textbf{Arbeitsblatt 4} (Theorie,
Plenumsboxen, Demoversuch, Dielektrikum-Herleitung). Grundprinzip: erst selbst
messen/beobachten, danach im Plenum vergleichen und die Formel festhalten --
nie umgekehrt.}

\vspace{0.3cm}

\renewcommand{\arraystretch}{1.3}
\setlength{\LTpre}{6pt}
\setlength{\LTpost}{6pt}

\begin{longtable}{@{}p{5.3cm}p{1.5cm}p{5.9cm}p{5.9cm}p{3.3cm}p{2.4cm}@{}}
\rowcolor{headgray}
\textbf{Inhalt} & \textbf{Zeit} & \textbf{Lehrperson} & \textbf{Schülerinnen und Schüler} & \textbf{Material} & \textbf{Sozialform} \\
\midrule
\endfirsthead

\rowcolor{headgray}
\textbf{Inhalt} & \textbf{Zeit} & \textbf{Lehrperson} & \textbf{Schülerinnen und Schüler} & \textbf{Material} & \textbf{Sozialform} \\
\midrule
\endhead

\bottomrule
\endfoot

\rowcolor{lightrow}
\textbf{Einstieg} \newline
\textit{Arbeitsblatt} -- Box \glqq Einstieg: Wozu ein Kondensator?\grqq{}
&
4'\newline{\footnotesize Ist:\dotfill}
&
Zeigt die Bilder (Defibrillator, Kondensatoren-Sortiment), fragt mündlich:
Wozu braucht es das neben einer Batterie? Was unterscheidet die
Kondensatoren im Bild?
&
Betrachten die Bilder, hören zu, äussern erste Vermutungen (kein
schriftlicher Auftrag).
&
Arbeitsblatt (Bilder)
&
Plenum
\\

\textbf{Theorie 1} \newline
\textit{Arbeitsblatt} -- Theoriebox \glqq Der Kondensator\grqq{}: Aufbau,
$U=E\cdot d$, Batterie-Skizzen (nicht angeschlossen/angeschlossen/getrennt).
&
8'\newline{\footnotesize Ist:\dotfill}
&
Erklärt den Plattenaufbau (±$Q$), zeigt $U=E\cdot d$ als bereits Bekanntes
(keine Herleitung), erklärt anhand der drei Skizzen Strom und Ladung beim
An-/Trennen der Batterie.
&
Hören zu, verfolgen die Skizzen im Arbeitsblatt. Keine eigene Aufgabe.
&
Arbeitsblatt
&
Lehrervortrag
\\

\rowcolor{lightrow}
\textbf{Lernaufgabe Teil 1} \newline
\textit{Aufgabenblatt} -- Intro-Box + \glqq Teil 1: Zusammenhang zwischen
$Q$, $U$ und $C$\grqq{}
&
10'\newline{\footnotesize Ist:\dotfill}
&
Zeigt Simulation und Excel-Blatt \texttt{Teil 1 - U und Q} zuerst am
Beamer, geht danach herum und unterstützt beim Messen.
&
Öffnen zu zweit Simulation und Excel-Datei. Messen $Q$ bei $\geq5$
$U$-Werten (Geometrie konstant), beschreiben den Zusammenhang \textbf{offen}
(noch keine vorgegebene Antwort), vergleichen $C$ mit der Simulationsanzeige.
&
Laptop, PhET-Simulation \glqq Capacitor Lab: Basics\grqq{}, Excel-Vorlage,
Aufgabenblatt
&
Partnerarbeit (PC)
\\

\textbf{Plenum: Kondensatoren im Vergleich} \newline
\textit{Arbeitsblatt} -- Groupbox, neues $Q$-$U$-Diagramm ($C_1$ vs. $C_2$)
&
6'\newline{\footnotesize Ist:\dotfill}
&
Moderiert den Vergleich am $Q$-$U$-Diagramm, lässt die Definitionsgleichung
$C:=Q/U$ (Einheit Farad) gemeinsam herleiten und notieren, verweist auf die
eigenen Messwerte aus Teil 1.
&
Erklären den Unterschied zwischen den beiden Kondensatoren im Diagramm,
notieren die Definitionsgleichung $C:=Q/U$ und die pF/nF/µF/mF-Präfixe.
&
Arbeitsblatt, Tafel/Beamer
&
Plenum
\\

\rowcolor{lightrow}
\textbf{Wassertankanalogie} \newline
\textit{Arbeitsblatt} -- Theoriebox (Bilder, Sortieraufgabe, Rollen-Frage)
&
8'\newline{\footnotesize Ist:\dotfill}
&
Leitet kurz in die Analogie ein (Wasserfass-Bild, Ventil offen/zu), lässt
die Sortieraufgabe zu zweit lösen, sammelt Antworten zur Rollen-Frage.
&
Betrachten die Analogie-Bilder, ordnen zu zweit Wassertank- und
Kondensator-Grössen zu (Sortieraufgabe), beantworten \glqq Welche Rolle
könnte ein Kondensator in einem Stromkreis spielen?\grqq
&
Arbeitsblatt
&
Partnerarbeit / Plenum
\\

\textbf{Theorie 2} \newline
\textit{Arbeitsblatt} -- Theoriebox \glqq Kapazität und Plattengeometrie:
die bekannten Bausteine\grqq{}
&
5'\newline{\footnotesize Ist:\dotfill}
&
Repetiert $E=Q/(\varepsilon_0\cdot A)$ und $U=E\cdot d$ als bereits
bekannt, kombiniert sie bewusst noch nicht, erklärt die Hinweisbox zum
homogenen Feld.
&
Hören zu, verfolgen die Repetition im Arbeitsblatt. Keine eigene Aufgabe.
&
Arbeitsblatt, Tafel/Beamer
&
Lehrervortrag
\\

\midrule
\multicolumn{6}{@{}c@{}}{\textit{--- Ende Lektion 1 (41') / Beginn Lektion 2 ---}} \\
\midrule

\rowcolor{lightrow}
\textbf{Lernaufgabe Teil 2} \newline
\textit{Aufgabenblatt} -- \glqq Teil 2: Einfluss von Fläche und
Abstand\grqq{}
&
10'\newline{\footnotesize Ist:\dotfill}
&
Hilft beim Wechsel zur nächsten Simulationsansicht und beim Blattwechsel
in Excel (\texttt{Teil 2a}, \texttt{Teil 2b}).
&
Messen zu zweit zuerst $C$ vs. $A$ (Excel-Blatt \texttt{Teil 2a - A und
C}), dann $C$ vs. $d$ (\texttt{Teil 2b - d und C}), beschreiben die
Zusammenhänge im Diagramm \textbf{offen} ($C\sim A$ bzw. $C\sim 1/d$
selbst erkennen).
&
Laptop, PhET-Simulation, Excel-Vorlage, Aufgabenblatt
&
Partnerarbeit (PC)
\\

\textbf{Plenum: Von der Geometrie zur Kondensatorformel} \newline
\textit{Arbeitsblatt} -- Groupbox, algebraische Herleitung
&
6'\newline{\footnotesize Ist:\dotfill}
&
Moderiert die Herleitung an der Tafel: die drei bekannten Formeln ($E$,
$U=E\cdot d$, $C=Q/U$) werden kombiniert, lässt Zwischenschritte
vorschlagen.
&
Kombinieren im Plenum die drei Formeln algebraisch zu
$C=\varepsilon_0\cdot A/d$, mit vollständigem Rechenweg.
&
Arbeitsblatt, Tafel/Beamer
&
Plenum
\\

\rowcolor{lightrow}
\textbf{Formel notieren} \newline
\textit{Aufgabenblatt} -- Teil 2, Marke \glqq Eure Kondensatorformel\grqq{}
&
2'\newline{\footnotesize Ist:\dotfill}
&
Kurzer Kontrollgang, ob alle die Formel richtig übertragen haben.
&
Tragen $C=\varepsilon_0\cdot A/d$ in eigener Handschrift ins Aufgabenblatt
ein.
&
Aufgabenblatt
&
Partnerarbeit (PC)
\\

\textbf{Lernaufgabe Teil 3} \newline
\textit{Aufgabenblatt} -- \glqq Teil 3: Kondensator von der Batterie
getrennt\grqq{}
&
6'\newline{\footnotesize Ist:\dotfill}
&
Erinnert bei Bedarf an die \glqq getrennt\grqq-Skizze aus Theorie 1
(Arbeitsblatt) als Gedankenstütze.
&
Laden den Kondensator wie in Teil 1, trennen ihn von der Batterie,
verändern $d$, begründen mit $Q=C\cdot U$, dass $Q$ konstant bleibt und
$U$ sich anpasst, prüfen es in der Simulation.
&
Laptop, PhET-Simulation, Aufgabenblatt
&
Partnerarbeit (PC)
\\

\rowcolor{lightrow}
\textbf{Experiment} \newline
\textit{Arbeitsblatt} -- Groupbox \glqq Experiment: Kondensator und
Lichtquelle\grqq{} (alle nach vorne)
&
6'\newline{\footnotesize Ist:\dotfill}
&
Führt den Demoversuch vorne durch (Kondensatoren laden, an Lichtquelle
anschliessen), holt zuerst die Vorhersagen ab (Predict-Observe-Explain).
&
Kommen nach vorne, sagen anhand der Wassertankanalogie voraus, ob die
Lichtquelle bei allen Kondensatoren gleich lange leuchtet, beobachten
danach den Versuch.
&
Kondensatoren (verschiedene Kapazitäten), Lichtquelle (LED),
Spannungsquelle, Schalter, Arbeitsblatt
&
Demo (Plenum) + Einzelvorhersage
\\

\textbf{Lernaufgabe Teil 4} \newline
\textit{Aufgabenblatt} -- \glqq Teil 4: Leuchtdauer im Simulationsscreen
`Light Bulb'\grqq{} \textit{(1. Streichkandidat bei Zeitdruck.)}
&
4'\newline{\footnotesize Ist:\dotfill}
&
Kurze Unterstützung beim Screen-Wechsel.
&
Wechseln zum Screen \glqq Light Bulb\grqq, laden zwei Kondensatoren
unterschiedlicher Kapazität bei derselben Spannung, trennen sie und lassen
sie über die Glühbirne entladen -- bestätigen die Demo-Beobachtung und
begründen sie zusätzlich mit $Q=C\cdot U$.
&
Laptop, PhET-Simulation, Aufgabenblatt
&
Partnerarbeit (PC)
\\

\rowcolor{lightrow}
\textbf{Theorie 3 + Anwendung} \newline
\textit{Arbeitsblatt} -- Theoriebox \glqq Das Dielektrikum\grqq{} +
Beispielbox \glqq Kapazitive Sensoren\grqq{} \textit{(2. Streichkandidat:
Anwendung nur kurz mündlich erwähnen.)}
&
8'\newline{\footnotesize Ist:\dotfill}
&
Erklärt anhand der Bilder die Polarisation und den Vergleich zum bereits
bekannten Faraday-Käfig (dort vollständige, hier nur teilweise
Feldauslöschung), danach kurz die Alltagsbeispiele (Kofferraumsensor,
Touchscreen).
&
Betrachten die Bilder (Polarisation, Wassermolekül als Dipol,
Faraday-Käfig-Vergleich, Kofferraumsensor/Touchscreen). Keine eigene
Aufgabe.
&
Arbeitsblatt, Tafel/Beamer
&
Lehrervortrag
\\

\textbf{Herleitung} \newline
\textit{Arbeitsblatt} -- Groupbox \glqq Herleitung (Partnerarbeit):
Kapazität mit Dielektrikum\grqq{}
&
9'\newline{\footnotesize Ist:\dotfill}
&
Geht herum, hilft insbesondere bei Teilaufgabe (a) (Skizze der
Nettoladung, mit Tipp-Hinweisbox) und (c) (Reaktion der Batterie).
&
Bearbeiten zu zweit: Nettoladung am Rand einzeichnen, Auswirkung auf $E$,
auf die freie Ladung, auf $C$, leiten zuletzt
$C=\varepsilon_r\cdot\varepsilon_0\cdot A/d$ über $\varepsilon_r:=C/C_0$
her.
&
Arbeitsblatt, Taschenrechner
&
Partnerarbeit
\\

\rowcolor{lightrow}
\multicolumn{1}{@{}p{5.3cm}}{\textbf{Zusammenfassung}} & \multicolumn{1}{p{1.5cm}}{\textbf{92'}} & \multicolumn{4}{@{}p{18.9cm}@{}}{\textit{Lektion 1: 41' (Phasen 1--6). Lektion 2: 51' (Phasen 7--14). Bei Zeitdruck zuerst Lernaufgabe Teil 4 streichen (bestätigt nur den unmittelbar vorher gezeigten Demoversuch), danach die Anwendungsbox nur kurz mündlich erwähnen statt im Detail zu besprechen.}} \\

\end{longtable}

\vspace{0.2cm}
{\small Mögliche Impulsfragen der Lehrperson und erwartbare Schülerfragen mit Antwort und Erklärung: siehe \texttt{kondensatoren-\allowbreak fragenkatalog.pdf}.}

\end{document}
